\PassOptionsToPackage{unicode=true}{hyperref} % options for packages loaded elsewhere
\PassOptionsToPackage{hyphens}{url}
%
\documentclass{book}        % from not too short p76 pagestyle
%\usepackage{fancyhdr}
%\pagestyle{fancy}
%ensure chapter section headngs in lower case
%
\usepackage{graphicx}
\DeclareGraphicsExtensions{.png,.jpg}
\usepackage{xeCJK}
\setCJKmainfont{SimSun}
\usepackage{lmodern}
\usepackage{amssymb,amsmath}
\usepackage{ifxetex,ifluatex}
\usepackage{fixltx2e} % provides \textsubscript
\ifnum 0\ifxetex 1\fi\ifluatex 1\fi=0 % if pdftex
  \usepackage[T1]{fontenc}
  \usepackage[utf8]{inputenc}
  \usepackage{textcomp} % provides euro and other symbols
\else % if luatex or xelatex
  \usepackage{unicode-math}
  \defaultfontfeatures{Ligatures=TeX,Scale=MatchLowercase}
\fi
% use upquote if available, for straight quotes in verbatim environments
\IfFileExists{upquote.sty}{\usepackage{upquote}}{}
% use microtype if available
\IfFileExists{microtype.sty}{%
\usepackage[]{microtype}
\UseMicrotypeSet[protrusion]{basicmath} % disable protrusion for tt fonts
}{}
\IfFileExists{parskip.sty}{%
\usepackage{parskip}
}{% else
\setlength{\parindent}{0pt}
\setlength{\parskip}{6pt plus 2pt minus 1pt}
}
\usepackage{hyperref}
\hypersetup{
            pdfborder={0 0 0},
            breaklinks=true}
\urlstyle{same}  % don't use monospace font for urls
\usepackage{longtable,booktabs}
% Fix footnotes in tables (requires footnote package)
\usepackage{multirow} %for tabular combine multirow
\IfFileExists{footnote.sty}{\usepackage{footnote}\makesavenoteenv{longtable}}{}
\setlength{\emergencystretch}{3em}  % prevent overfull lines
\providecommand{\tightlist}{%
  \setlength{\itemsep}{0pt}\setlength{\parskip}{0pt}}
\setcounter{secnumdepth}{0}
% Redefines (sub)paragraphs to behave more like sections
\ifx\paragraph\undefined\else
\let\oldparagraph\paragraph
\renewcommand{\paragraph}[1]{\oldparagraph{#1}\mbox{}}
\fi
\ifx\subparagraph\undefined\else
\let\oldsubparagraph\subparagraph
\renewcommand{\subparagraph}[1]{\oldsubparagraph{#1}\mbox{}}
\fi

% set default figure placement to htbp
\makeatletter
\def\fps@figure{htbp}
\makeatother
\date{}

\usepackage{graphicx}


\usepackage{amsmath}


\raggedbottom

\begin{document}               % plus the \end{document} command at the end.

与成为有效过渡者的要素一样，企业、团队、个人要有效改善，必须有针对性，大部分企业没有利用数据、找根因、做改进。所以开始时应只针对缺陷排除和返工工时，每次迭代回顾时分析根因。

``听你这样说，反观我们团队，我立马想到有很多地方需要改进，例如需求调研、开发流程、持续集成、测试自动化、配置管理等。为什么你提议只针对减少缺陷引起的返工做改进？''

要理解背后的原因，先听听质量大师Dr Juran回顾过去质量改善的成功要素：

\begin{itemize}
\tightlist
\item
  All Improvement takes place Project by Project
  所有改进只能按项目逐个进行
\item
  Backlog of Improvement projects is huge 潜在的改进项目非常多
\item
  Quality Improvement does not come Free 质量改进并非免费
\item
  Major gains come from the Vital few projects
  主要效果（收益）来自少数重要项目
\item
  Quality improvement extends to All parameters 质量改进可扩展到各类参数
  （如生产率、进度、安全性、环境等）
\end{itemize}

（以上源自 Juran 博士《 Juran Quality Handbook 》第五章质量改进。）

所以如果开始时同时对4 -- 5
个方向同时做改进，反而会失去聚焦，难以尽早看到改进的效果。

但如果只针对降低缺陷引起的返工，目标明确不仅让大家容易接受，而且也容易扩展到各个过程，例如需求、开放、产品集成，甚至配置管理。

因开始时能否引起高层对改进的兴趣很重要，与企业领导的首次见面，我也都会讲以下内容，说明团队如何分析迭代缺陷数据，开始定量管理。\\
Ct (\(= Ct/2 + Ct/2\))


\framebox{%
\begin{minipage}[t]{0.97\columnwidth}\raggedright
\(\cos \theta =  \cos \varphi _A \cos \varphi _B \cos\)(
\(\lambda _A\) - \(\lambda _B\) ) +
\(\sin \varphi _A \sin \varphi _B\)

\begin{description}
\item[]
\begin{description}
\tightlist
\item[]
\quad\,=\(\cos 39.905556^\circ \times  \cos48.8583^\circ \times \cos\)(\(116.424722^\circ  - 2.29451^\circ\)
) \(+ \sin 39.905556^\circ \times  \sin48.8583 ^\circ\)
\end{description}
\end{description}

\[= 0.767103 \times 0.657923 \times\]（\(-0.40881\)）\(+ 0.64152 \times 0.758084\)

\[= 0.28\]

\[\theta = 1.287\] (rad) AB 距离 = 地球的半径
\(6,378KM \times 1.287 = 8,208KM\) （实际距离是 8,217KM）
\strut
\end{minipage}}

\begin{equation}
\begin{split}
    V_{i,t+1} & =   w * V_{i,t} + c_1 * r_1 * (Pbest_{i,t} - X_{i,t}) \\
    & + c_2 * r_2 * (Gbest_t - X_{i,t})
\end{split}
\end{equation}




\end{document}
